\documentclass[11pt,a4paper]{article}
\usepackage[utf8]{inputenc}
\usepackage[T1]{fontenc}
\usepackage{lmodern}
\usepackage{amsmath,amssymb}
\usepackage{geometry}
\geometry{margin=1in}
\usepackage{xcolor}
\usepackage{hyperref}
\usepackage{cleveref}
\usepackage{booktabs}
\usepackage{array}
\usepackage{microtype}
\usepackage{tcolorbox}

\hypersetup{colorlinks=true, linkcolor=blue, citecolor=red, urlcolor=cyan,
  pdftitle={Calibration Checks and a Literature Recount for a Mixed Fuzzy-Dark-Matter Window},
  pdfauthor={Xavier Callens and the SocrateAI Scientific Agora Collaboration}}

\newtcolorbox{scopebox}{colback=gray!5, colframe=gray!45, boxrule=0.4pt, arc=1pt,
  left=6pt, right=6pt, top=3pt, bottom=3pt, fonttitle=\bfseries\small, title={What is not claimed}}
\newtcolorbox{headlinebox}{colback=red!4, colframe=red!45!black, boxrule=0.5pt, arc=1pt,
  left=6pt, right=6pt, top=4pt, bottom=4pt, fonttitle=\bfseries\small, title={Headline result}}
\newtcolorbox{correctionbox}{colback=orange!6, colframe=orange!55!black, boxrule=0.5pt, arc=1pt,
  left=6pt, right=6pt, top=4pt, bottom=4pt, fonttitle=\bfseries\small, title={Corrections recorded in band}}

\title{\Large \textbf{Calibration Checks and a Literature Recount}\\[0.2em]
\textbf{for a Mixed Fuzzy-Dark-Matter Search Window}\\[0.3em]\large Stream 3 status note}
\author{\textbf{Xavier Callens} \and \textbf{SocrateAI Scientific Agora Collaboration} \\[0.3em]
\textit{Stream 3 --- Experimentation (\texttt{DarkMatterK3-Home.github.io})}}
\date{2026-10-08 \quad (preprint; not reviewed)}

\begin{document}
\maketitle

\begin{scopebox}
This note contains no dark-matter result. Nothing in it compares a model with data: no prediction is pinned
for the search discussed here, no output carries a TEST or FIT label, and the comparison phases of that
search are not authorized. The calibration checks are model checks of an exactly solvable Hamiltonian, not
evidence about any physical system. The $K3$ surfaces selected by the geometry stream, including $X_3$ and
$X_4$, map to no dark-matter observable, and no number from that work is used here as a prior, a parameter
or an input.
\end{scopebox}

\begin{abstract}
We report the state of the programme's experimentation stream on 2026-10-08. Two numerical calibrations
pass on finite Fibonacci rings of the Aubry--Andr\'e chain. The exact duality between couplings $\lambda$ and
$4J^2/\lambda$ holds with errors below $10^{-11}$ at $N\le233$, and a seeded random potential of the same strength
breaks it, as required. The extended and localized regimes and their crossing at $\lambda=2J$ hold for
$N\le987$ outside a band $|\lambda-2J|<0.15$, which is reported and not claimed. Separately, a literature
recount was carried out for a proposed search for a sub-dominant fuzzy-dark-matter component with fraction
$f<1$ in the one-dimensional Lyman-$\alpha$ flux power spectrum. Published mixed-fraction bounds reduce the
number of grid cells that would be both decisive and still open from $221$ to $182$ using text-stated bounds,
or to $152$ when figure-read contours are included. The proposal's stop condition does not fire. No published
analysis of the DESI DR1 spectrum for a mixed fraction was found.
\end{abstract}

\section{Where the stream stands}\label{sec:status}

The stream's governing rule is that no real-data comparison is coded before a prediction is pinned. Two of
its routes are stopped: a two-dimensional transverse search is closed by the data floors of public surveys,
and a flux/tadpole route is blocked until a threefold base is specified. The open direction is an exclusion-type
search over a grid of fuzzy-dark-matter masses $m$ and fractions $f$. Under its proposal, Phase~0 is a
literature re-survey; Phases~1--4 build and run the comparison. The decision authority approved Phase~0 only,
on 2026-10-08. A noise-calibration experiment (E2) is held because no open noise dataset exists, and two
laboratory-analogue experiments (E3, E4) remain appendices.

\section{Calibration of the Aubry--Andr\'e arm}\label{sec:e1}

The model is the tight-binding chain $(H\psi)_n = J(\psi_{n+1}+\psi_{n-1}) + \lambda\cos(2\pi pn/N)\,\psi_n$ on a
ring of $N$ sites, with $(N,p)$ consecutive Fibonacci numbers so that $\gcd(p,N)=1$. The discrete Fourier
transform intertwines $H(J,\lambda)$ with $H(\lambda/2,2J)$ exactly, so the spectra at $\lambda$ and $4J^2/\lambda$
agree after rescaling, and at $\lambda=2J$ each non-degenerate eigenvector has equal inverse participation
ratios in position and momentum. These are algebraic identities; the checker verifies them in floating point,
so every verdict is PASS($N$) at a stated tolerance. Negative controls run first and must fail: a seeded
random potential of matched strength breaks the duality, and a non-coprime $p$ makes the checker refuse to
run.

\begin{table}[h]
\centering\small
\begin{tabular}{rrlllll}
\toprule
$N$ & $p$ & max duality gap & max intertwiner gap & control gap (random) & controls fail? & verdict \\
\midrule
$13$ & $8$ & $9.0\times10^{-15}$ & $2.9\times10^{-14}$ & $2.0\times10^{-1}$ & yes & PASS \\
$89$ & $55$ & $3.9\times10^{-14}$ & $1.0\times10^{-12}$ & $2.8\times10^{-1}$ & yes & PASS \\
$233$ & $144$ & $2.3\times10^{-13}$ & $7.8\times10^{-12}$ & $4.7\times10^{-1}$ & yes & PASS \\
\bottomrule
\end{tabular}

\caption{E1, tolerance $10^{-9}$. Gaps are maxima over $\lambda\in\{0.5,1,3\}J$; the control column is the
smallest duality gap produced by the random potential, which must be large.}\label{tab:e1}
\end{table}

E1b extends the check to the whole coupling range and to the localization dichotomy, with acceptance criteria
written in the checker before it ran. Extended regime: the product $N\cdot$(mean IPR) varies by less than a
factor $1.5$ across rings. Localized regime: the mean IPR is stable within $10\%$ between the two largest rings
and exceeds $0.1$. The localized regime is evaluated at a generic phase, because at phase zero exactly
degenerate mirror pairs make the eigenbasis, and therefore the mean IPR, ill-defined.

\begin{table}[h]
\centering\small
\begin{tabular}{rlllp{4.2cm}}
\toprule
$\lambda/J$ & regime & statistic over $N\in\{89,233,377,987\}$ & criterion met & status \\
\midrule
$0.5$ & extended & $\max/\min$ of $N\cdot\mathrm{IPR}$ $=1.024$ & yes & claimed \\
$1.0$ & extended & $\max/\min$ of $N\cdot\mathrm{IPR}$ $=1.038$ & yes & claimed \\
$1.5$ & extended & $\max/\min$ of $N\cdot\mathrm{IPR}$ $=1.056$ & yes & claimed \\
$1.9$ & extended & $\max/\min$ of $N\cdot\mathrm{IPR}$ $=1.318$ & yes & near-critical: reported, not claimed \\
$2.1$ & localized & $\mathrm{IPR}\in[0.149,0.191]$ & no & near-critical: reported, not claimed \\
$2.5$ & localized & $\mathrm{IPR}\in[0.390,0.409]$ & yes & claimed \\
$3.0$ & localized & $\mathrm{IPR}\in[0.528,0.542]$ & yes & claimed \\
$4.0$ & localized & $\mathrm{IPR}\in[0.669,0.680]$ & yes & claimed \\
\bottomrule
\end{tabular}

\caption{E1b scaling ($J=1$). The crossing of mean IPR in position and momentum changes sign at $\lambda=2J$
for every coupling tested. Rows in the near-critical band are reported, not claimed.}\label{tab:e1b}
\end{table}

These checks calibrate numerical machinery. They are not a measurement, they say nothing about any $K3$
surface, and no result here is evidence for a dark-matter model.

\section{Phase 0: a literature recount for mixed fuzzy dark matter}\label{sec:p0}

The proposed search covers $13$ masses from $10^{-22}$ to $10^{-19}$\,eV and $20$ fractions $f$ from $0.05$ to
$1$. A cell is \emph{decisive} if an earlier pre-flight, using an optimistic proxy, finds that the data could
reach two standard deviations there. A cell is \emph{open} if no published constraint already excludes it. For
$f=1$ the window is already covered by published bounds; the question for Phase~0 was how much of the $f<1$
region published mixed-fraction analyses have closed since the pre-flight table was written.

We re-read the sources and tabulated each bound with its locator and quoted sentence. Text-stated bounds are
kept apart from contours read off figures. The latter come from a colour-thresholding script run on
hash-pinned PDFs, with a reading error of about $0.03$ in $f$. Every figure value is shifted toward closing
fewer cells, and the script refuses to run unless its reading of one contour reproduces the two limits that
paper states in its text. The main sources are Kobayashi et al.\ (2017), for which a fuzzy component above $30\%$
requires $m\gtrsim10^{-21}$\,eV, and Liu, Gong and Zhou (2026), whose $95\%$ contour gives no effective limit
above $10^{-20}$\,eV.

\begin{table}[h]
\centering\small
\begin{tabular}{lrl}
\toprule
Count & cells & P0 fires? \\
\midrule
decisive cells (optimistic proxy) & $258$ & --- \\
decisive and open, before Phase 0 & $221$ & --- \\
decisive and open, text-grade bounds & $182$ & no \\
decisive and open, text and figure reads & $152$ & no \\
text-grade, frequentist limits instead of Bayesian & $188$ & --- \\
\bottomrule
\end{tabular}

\caption{Phase~0 recount. ``Decisive'' is the pre-flight's optimistic proxy, not a forecast. The stop
condition fires if fewer than $10$ cells survive or if all survivors have $f\ge0.5$.}\label{tab:p0}
\end{table}

\begin{headlinebox}
Phase~0 leaves $182$ decisive and open cells using text-stated bounds, or $152$ with figure reads. Survivors
reach down to $f=0.05$ at most masses, so the stop condition does not fire. The DESI DR1 cosmological analysis
of the one-dimensional flux power spectrum states that constraining the nature of dark matter is beyond its
scope, so the proposal's secondary trigger does not fire either.
\end{headlinebox}

What these counts are worth: four parallel survey reports covering CMB and large-scale structure, dwarf
galaxies, reviews and the Lyman-$\alpha$ forest were not all audited against their sources, and nothing from
them is applied. An unaudited bound, if real, could only close more cells, so the counts are upper bounds on
how open the grid is. An Eridanus~II star-cluster constraint is recorded and not applied, because its figure
carries resonance bands and an approximation caveat near a few $10^{-21}$\,eV.

\section{What happens next, and what does not}\label{sec:next}

Phases~1--4 have not started and are not authorized: they need their predecessors' artifacts, a pinned
prediction amendment, and the decision authority's answers to five open questions of the proposal. Any output
of those phases would be labeled as an exclusion or FIT, never as a TEST, until the pin. E2 stays held. The
geometry stream's results of the same date (\texttt{papers/stream2\_k3\_identification\_twisted\_route\_2026\_10\_08.pdf}
in the geometry repository) have their certificates mirrored under \nolinkurl{data/mirrors/stream2/} for the
record; they are not inputs to anything above.

\begin{correctionbox}
\begin{enumerate}
\item E1b, first run: the localized-regime criterion failed at phase zero. The cause was eigenbasis
  dependence inside exactly degenerate mirror pairs, not a physics failure. The criterion was then evaluated at
  a generic phase; the phase-zero values are kept in the certificate.
\item E1b, second run: at $\lambda=2.1J$ the localized criterion still fails on rings up to $987$ sites. The band
  $|\lambda-2J|<0.15$ is reported and not claimed; the criterion itself was not changed.
\item Phase~0 report draft: a sentence saying the two Lyman-$\alpha$ contours agreed below $10^{-21}$\,eV was
  false (at $10^{-21}$\,eV they read about $0.35$ and $0.63$). It was corrected before the report was merged.
\item Phase~0 figure script: the pinned hash of one PDF was first typed with a missing digit; the script's own
  hash check refused to run, which is the purpose of the check.
\item The repository's continuous-integration job fails at its dependency-installation step on every recent
  run, before any test executes. A fix is prepared and waits for a token permission. Local runs of the full
  test suite pass; the CI job is not green and is not described as green.
\end{enumerate}
\end{correctionbox}

\section*{Reproducibility}
\begin{itemize}\small
\item Tables: \nolinkurl{scripts/render_stream3_note_tables.py} (\texttt{--check} fails on drift); controls in
  \nolinkurl{checkers/tests/test_render_stream3_note_tables.py}.
\item Calibrations: \nolinkurl{checkers/check_aubry_andre_selfduality.py},
  \nolinkurl{checkers/check_aubry_andre_localization_scaling.py}.
\item Phase~0: \nolinkurl{scripts/wp_e6_v2_p0_overlay.py}, \nolinkurl{scripts/wp_e6_v2_p0_digitize_figures.py},
  report \nolinkurl{briefs/WP_E6_V2_PHASE0_RESURVEY_REPORT_2026_10_08.md}.
\end{itemize}

\section*{Provenance}
Generated-by: Claude (Opus~5.5), Stream~3 support \textbar\ Verified-by: the renderer's \texttt{--check} and
controls; each cited checker's negative controls \textbar\ Reviewed-by: N (preprint; not reviewed).

\begin{thebibliography}{9}
\bibitem{AA} S.~Aubry and G.~Andr\'e, ``Analyticity breaking and Anderson localization in incommensurate
  lattices,'' \textit{Ann. Israel Phys. Soc.} \textbf{3} (1980) 133.
\bibitem{Kob} T.~Kobayashi, R.~Murgia, A.~De~Simone, V.~Ir\v{s}i\v{c}, M.~Viel, ``Lyman-alpha constraints on
  ultralight scalar dark matter: implications for the early and late universe,'' \textit{arXiv:1708.00015}.
\bibitem{Liu} Liu, Gong and Zhou, mixed fuzzy-dark-matter constraints from the high-redshift Lyman-$\alpha$
  forest, \textit{arXiv:2606.06969}.
\bibitem{DESI} J.~Chaves-Montero et al.\ (DESI), ``Cosmological analysis of the DESI DR1 Ly$\alpha$ 1D power
  spectrum,'' \textit{arXiv:2601.21432}.
\bibitem{DESIm} N.~G.~Kara\c{c}ayl\i{} et al.\ (DESI), ``DESI DR1 Ly$\alpha$ 1D power spectrum: the optimal
  estimator measurement,'' \textit{arXiv:2505.07974}.
\bibitem{MN} D.~J.~E.~Marsh and J.~C.~Niemeyer, ``Strong constraints on fuzzy dark matter from ultrafaint
  dwarf galaxy Eridanus II,'' \textit{arXiv:1810.08543}.
\end{thebibliography}

\end{document}
